\ticket{5}{Последовательная и параллельная сложность алгоритмов}

\begin{examplan}
    \item Определить последовательную сложность и асимптотические оценки.
    \item Ввести параллельное время, работу, стоимость, ускорение и эффективность.
    \item Объяснить информационный граф алгоритма и критический путь.
    \item Раскрыть ресурс параллелизма и его связь с достижимым ускорением.
\end{examplan}

\subsection*{Короткая версия для письменного ответа}

\textbf{Последовательная сложность} оценивает ресурсы алгоритма при выполнении на одном процессоре: время $T(n)$ и память $S(n)$ как функции размера входа $n$.

\textbf{Параллельная сложность} оценивает выполнение алгоритма на нескольких процессорах с учетом зависимостей, коммуникаций и синхронизации. Главные величины: работа $T_1(n)$, параллельное время $T_p(n)$ и критический путь $T_\infty(n)$.

\subsection*{Последовательная сложность}

Для сравнения алгоритмов используют асимптотический анализ:
\begin{itemize}
    \item $O(g(n))$ --- верхняя граница роста;
    \item $\Omega(g(n))$ --- нижняя граница;
    \item $\Theta(g(n))$ --- точная асимптотическая граница.
\end{itemize}

Типичные классы: $O(1)$, $O(\log n)$, $O(n)$, $O(n\log n)$, $O(n^2)$, полиномиальные, экспоненциальные и факториальные. Полиномиальные алгоритмы обычно считаются практически более приемлемыми, чем экспоненциальные, хотя реальная применимость зависит от констант и размера входа.

\subsection*{Параллельная сложность}

\begin{description}
    \item[$T_p(n)$] время выполнения на $p$ процессорах.
    \item[$T_1(n)$] работа, то есть суммарное число операций; это время на одном процессоре для данного алгоритма.
    \item[$C_p(n)$] стоимость: $C_p(n)=pT_p(n)$.
    \item[$S_p(n)$] ускорение: $S_p(n)=T_1(n)/T_p(n)$ или отношение лучшего последовательного времени к параллельному.
    \item[$E_p(n)$] эффективность: $E_p(n)=S_p(n)/p$.
\end{description}

Идеальное ускорение равно $p$, но на практике оно ограничено последовательными участками, зависимостями данных, синхронизацией и обменами.

\textbf{Закон Амдала.} Если доля нераспараллеливаемой части равна $f$, то максимальное ускорение при бесконечном числе процессоров ограничено:
\[
S_{\max} \leq \frac{1}{f}.
\]

\subsection*{Информационный граф}

\textbf{Информационный граф алгоритма} --- ориентированный ациклический граф, где вершины соответствуют элементарным операциям, а дуги --- информационным зависимостям. Дуга $u \to v$ означает, что операция $v$ не может начаться до получения результата операции $u$.

Информационный граф позволяет определить:
\begin{itemize}
    \item какие операции можно выполнять параллельно;
    \item общую работу $T_1$ как сумму весов вершин;
    \item критический путь $T_\infty$ как самый длинный путь по зависимостям;
    \item нижнюю оценку параллельного времени даже при неограниченном числе процессоров.
\end{itemize}

Пример: в выражении $(a+b)(c+d)$ сложения $a+b$ и $c+d$ независимы и могут выполняться параллельно, а умножение зависит от обоих результатов.

\subsection*{Ресурс параллелизма}

\textbf{Ресурс параллелизма} показывает среднее количество операций, которое потенциально можно выполнять одновременно:
\[
\Pi(n)=\frac{T_1(n)}{T_\infty(n)}.
\]

Он задает верхнюю границу ускорения:
\[
S_p(n) \leq \min(p,\Pi(n)).
\]

Если процессоров больше, чем ресурс параллелизма, часть процессоров будет простаивать.

Для суммирования $n$ чисел деревом: $T_1(n)=O(n)$, $T_\infty(n)=O(\log n)$, поэтому $\Pi(n)=O(n/\log n)$.

\begin{important}
    \item Последовательная сложность говорит о работе на одном процессоре, параллельная --- о времени с учетом зависимостей.
    \item Работа $T_1$ и критический путь $T_\infty$ --- две главные характеристики параллельного алгоритма.
    \item Информационный граф показывает зависимости между операциями.
    \item Критический путь нельзя ускорить добавлением процессоров.
    \item Ресурс параллелизма ограничивает максимально полезное число процессоров.
\end{important}

